% Samuel D. V. Stephens's Curriculum Vitae --- Cormorant Garamond edition
% Adapted from Geoff Boeing's template (https://github.com/gboeing/cv).
%
% IMPORTANT: this variant must be compiled with xelatex or lualatex
% (NOT pdflatex), because it loads Cormorant Garamond via fontspec.
% Cormorant Garamond is the same font used by the older prometheus_cv
% template in ~/Downloads/hnst/main.tex; it is a free font from
% https://github.com/CatharsisFonts/Cormorant
%
% Build:  python3 scripts/build-cv.py   (PDF and HTML; see the README)
%
\RequirePackage[l2tabu,orthodox]{nag} % warn if using any obsolete commands
\documentclass[11pt,letterpaper]{report} % document style
\usepackage{inputenc}
\usepackage{siunitx}							% package to properly set units
\usepackage{fontawesome5}						% package for icons (see list of available icons here: http://mirrors.ibiblio.org/CTAN/fonts/fontawesome5/doc/fontawesome5.pdf)
\usepackage[super]{nth}		
% fontspec replaces fontenc / inputenc under xelatex / lualatex
\usepackage{fontspec}
% Fonts are loaded by file name, so xelatex finds them in the TeX Live tree
% (fonts/truetype/catharsis/cormorantgaramond) on any machine, without
% installing them system-wide.
\setmainfont{CormorantGaramond-Regular.ttf}[
    BoldFont=CormorantGaramond-Bold.ttf,
    ItalicFont=CormorantGaramond-Italic.ttf,
    BoldItalicFont=CormorantGaramond-BoldItalic.ttf]
% give \sffamily (used by the Boeing section/subsection headers) a Garamond
% face so headers stay in the Cormorant family rather than falling back to
% a default sans face
\setsansfont{CormorantGaramond-Regular.ttf}[
    BoldFont=CormorantGaramond-Bold.ttf,
    ItalicFont=CormorantGaramond-Italic.ttf,
    BoldItalicFont=CormorantGaramond-BoldItalic.ttf]
\setmonofont{lmmono10-regular.otf}[BoldFont=lmmonolt10-bold.otf, ItalicFont=lmmono10-italic.otf]
\newfontfamily\GaramondLight{CormorantGaramond-Light.ttf}
\newcommand{\textlf}[1]{{\GaramondLight #1}}

% load babel, csquotes, and microtype in order
\usepackage[USenglish]{babel} % auto-regionalize hyphens, quote marks, etc
\usepackage[strict,autostyle]{csquotes} % smart and nestable quote marks
% microtype works with xelatex but only protrusion is supported; expansion
% and font-tracking via \SetTracking are not. We still load it for protrusion.
\usepackage[protrusion=true,expansion=false]{microtype}

% load everything else
\usepackage{datetime} % enable formatting of date output
\usepackage{enumitem} % enumerate with [resume] option
\usepackage{geometry} % manually set page margins
\usepackage{setspace} % custom line spacing
\usepackage{tabto}    % make nice tabbing
\usepackage{titlesec} % allow custom section fonts

% load hyperref last for compatibility
\usepackage{hyperref} % enable hyperlinks and pdf metadata

% what is your name?
\newcommand{\myname}{Samuel D. V. Stephens}

% how far to tab for list items with left-aligned date; Cormorant Garamond
% renders these dates a touch wider than EB Garamond, so give it room
\newcommand{\listtabwidth}{1.9cm}

% define font to use as document's title
\newcommand{\namefont}[1]{{\normalfont\bfseries\Huge{#1}}}

% Letter-spaced headers under xelatex: microtype's \lsstyle / \SetTracking
% are unsupported, so we define a Garamond face with an actual OpenType
% letter-spacing feature and use it in the heading definitions below.
\newfontfamily\headingfont{CormorantGaramond-Bold.ttf}[
    LetterSpace=8.0,
    WordSpace=1.4]

% set section heading fonts and before/after spacing
\titleformat{\section}{\headingfont\small\uppercase}{}{}{}{}
\titlespacing{\section}{0pt}{30pt plus 4pt minus 4pt}{8pt plus 2pt minus 2pt}

% set subsection heading fonts and before/after spacing
\titleformat{\subsection}{\headingfont\footnotesize}{}{}{}{}
\titlespacing{\subsection}{0pt}{16pt plus 4pt minus 4pt}{4pt plus 2pt minus 2pt}

% set page margins (assumes letter paper)
\geometry{body={6.5in, 9.0in},
    left=1.0in,
    top=1.0in}

% prevent paragraph indentation
\setlength\parindent{0em}

% set line spacing
\setstretch{0.9}

% define space between list items
\newcommand{\listitemspace}{0.25em}

% make unordered lists without bullets and use compact spacing
\renewenvironment{itemize}
{\begin{list}{}{\setlength{\leftmargin}{0em}
                \setlength{\parskip}{0em}
                \setlength{\itemsep}{\listitemspace}
                \setlength{\parsep}{\listitemspace}}}
{\end{list}}

% make tabbed lists so content is left-aligned next to years
\TabPositions{\listtabwidth}
\newlist{tablist}{description}{3}
\setlist[tablist]{leftmargin=\listtabwidth,
    labelindent=0em,
    topsep=0em,
    partopsep=0em,
    itemsep=\listitemspace,
    parsep=\listitemspace,
    font=\normalfont}

% placeholder command for sections without content yet
\newcommand{\forthcoming}{\textit{Forthcoming.}}

% print only the month and year when using \today
\newdateformat{monthyeardate}{\monthname[\THEMONTH] \THEYEAR}

% define hyperlink appearance and metadata for pdf properties
\hypersetup{
    colorlinks  = true,
    urlcolor    = black,
    citecolor   = black,
    linkcolor   = black,
    pdfauthor   = {Samuel Desiderius Vichaama Stephens},
    pdfkeywords = {neurobioengineering, mathematical physics, philosophy, cognitive science, biostatistics, education policy, polyadic algebra},
    pdftitle    = {\myname: Curriculum Vitae},
    pdfsubject  = {Curriculum Vitae},
    pdfpagemode = UseNone
}

\begin{document}
    \raggedright{}


    % display name as the document title
    \namefont{\myname}

    % affiliation and contact info blocks
    \vspace{1em}
    \begin{minipage}[t]{0.700\textwidth}
        % current primary affiliation, left-aligned
        Department of Neurology \\
        School of Medicine \\
        Stanford University
    \end{minipage}
    \begin{minipage}[t]{0.295\textwidth}
        % contact info details, right-aligned
        \raggedleft{}
        \href{mailto:samuelstephens2004@gmail.com}{samuelstephens2004@gmail.com} \\
        +1 214 843 5006 \\
        \href{https://www.samuel-stephens.com}{samuel-stephens.com}
    \end{minipage}

    \section*{Education}
    \begin{tablist}
        \item[2026--] \tab{}\textbf{B. S. in Mathematics and Music}, Salem State University, Salem, Massachusetts. Current GPA:~3.94. Coursework: Graduate Measure Theory (directed study, Folland's \textit{Real Analysis}); Philosophy of Politics and Society; Neuropsychology; American Romanticism; Piano Performance.
        \item[2019--22] \tab{}\textbf{High School Diploma, Pre-Law Track}, Judge Barefoot Sanders Law Magnet, Dallas, Texas. Graduated with Honors, Distinction, and a Distinguished Level of Achievement per the Texas Education Agency. Coursework included dual credit and 13 AP exams.
    \end{tablist}

    \section*{Appointments and Research Experience}
    %\subsection*{Academic Appointments}
    %\forthcoming{}

    \subsection*{Research Positions}
    \begin{tablist}
        \item[2026--]
        \tab{}\small{\textbf{\textsc{Mass General Brigham}}} \\[0.15em]
                             \small{\textit{Assistant Research Intern, Department of Neurology}} \\
                             Large-scale real-data phase of the CICADAS project (CICADAS-III), using MGB computing resources for large ICU EEG and clinical datasets. The framework was first run on real patient records in October 2026.
        \item[2026--]   \tab{}\small{\textbf{\textsc{Stanford University}}} \\[0.15em]
                              \small{\textit{Assistant Researcher, Department of Neurology, School of Medicine}} \\
                              Continuing the CICADAS-II line of work begun at BIDMC/HMS toward real-world clinical use; introduced personalized adaptive thresholds based on SOFA scores and age, and advancing implementation against EEG and the MORGOTH large-scale ICU epileptic dataset. The joint methods paper that grounds this work is under major revision at the \textit{Journal of Neural Engineering}.
        \item[2025--26] \tab{}\small{\textbf{\textsc{Beth Israel Deaconess Medical Center \& Harvard Medical School}}} \\[0.15em]
                              \small{\textit{Assistant Researcher, Department of Neurology}}, under Dr.~M.~Brandon Westover \\
                              Developed CICADAS-II, the successor to CICADAS (Causal Inference for Critical-Care Anti-Seizure Treatment with Disease Dynamics, Automated PKPD/Trial Simulation, and Survival Optimization): methodological frameworks that recover unbiased causal treatment effects from observational ICU data, with PKPD-aware modeling for diseased patients.
    \end{tablist}

    \subsection*{Paused or Discontinued Research}
    \begin{tablist}
        \item[2022--24] \tab{}\small{\textbf{\textsc{The University of Texas at Dallas}}}\\[0.15em]
                              \textit{Unpaid Research and Literature Review Collaborator}, under Dr.~Abdi \\
                              Meta-analysis on the Bouba-Kiki effect in children with Misophonia and/or Autism Spectrum Disorder.
    \end{tablist}

    \subsection*{Pre-College Research and Internships}
    \begin{tablist}
        \item[2021--22] \tab{}\small{\textbf{\textsc{School for Health Professions, Dallas}}} \\[0.15em]
                              \textit{Student Researcher}, Microbiology Department, under Dr.~Harrington (high-school research) \\
                              Identified novel functions of bacteria in the gut microbiome of \textit{Z.\ atratus} relevant to low-density polyethylene and polystyrene biodegradation. Written up as the AP Research academic paper submitted to the College Board.
        \item[2021--22] \tab{}\small{\textbf{\textsc{City of Dallas}}} \\[0.15em]
                              \textit{Dallas City Council Research Intern}, District 11 (high-school internship program through Judge Barefoot Sanders Law Magnet) \\
                              Researched smart-city development initiatives under the supervision of the city manager and the District 11 council member
    \end{tablist}

    \section*{Additional Education}
    \begin{tablist}
        \item[2026]     \tab{}\small{\textbf{\textsc{Massachusetts Institute of Technology}}} \\
                              Coursework in Bioengineering (Audit) \\
                              \textit{20.205[J]/9.26: Principles and Applications of Genetic Engineering for Biotechnology and Neuroscience}
        \item[2025--26] \tab{}\small{\textbf{\textsc{Harvard University}}}\\
                              Coursework in Mathematics and History (Audit) \\
                              \textit{Studies in Algebra and Group Theory (Math 55a); Studies in Real and Complex Analysis (Math 55b); Critical Theory (History 89)}
        \item[2025]     \tab{}\small{\textbf{\textsc{Wabash College}}} \\
                              Coursework in Mathematics, Neurobiology, and Music \\
                              \textit{Complex Analysis; Non-Fiction Creative Writing; Linguistics; Sociolinguistics; Principles of Neuroscience; Organ Performance}
        \item[2024]     \tab{}\small{\textbf{\textsc{Institute for Social Ecology}}} \\
                              Coursework in Sociology and Social Ecology \\
                              \textit{Legacies of Environmental Radicalism}
        \item[2024]     \tab{}\small{\textbf{\textsc{The University of Texas at Dallas}}} \\
                              Coursework in Mathematics (Audit) \\
                              \textit{Combinatorics; Advanced Calculus}
        \item[2023--24] \tab{}\small{\textbf{\textsc{Kenyon College}}} \\
                              Coursework in Mathematics, Philosophy, Political Science, and Music \\
                              \textit{Abstract Algebra I; Contemporary Mathematics Seminar; Symbolic Logic; Philosophy of Natural Science; Piano Practicum; Violin Practicum; Quest for Justice; STEM Research Methods}
    \end{tablist}

    \section*{Shared Course Notes}
    Personal notes on courses taken, audited, or followed out of interest, shared openly at \href{https://www.samuel-stephens.com/notes}{samuel-stephens.com/notes}. Mostly transcriptions and summaries of other people's lectures and texts; course numbering inspired by MIT.
    \begin{tablist}
        \item[2025--26] \tab{}\textit{Math 55a: Studies in Algebra and Group Theory}: the textbook listed under Books, with transcriptions of Professor Denis Auroux's lecture notes and review sheets (full-text PDF also in the \href{https://github.com/sdvstephens/uni}{\texttt{uni}} repository)
        \item[2026]     \tab{}\textit{Math 55b: Studies in Real and Complex Analysis}: transcribed lecture notes
        \item[2026]     \tab{}\textit{Critical Theory}: 25 lecture write-ups on the Frankfurt School (Luk\'acs, Horkheimer, Benjamin, Adorno, Marcuse, Habermas, and beyond)
        \item[2026]     \tab{}\textit{Basic Neuroscience and Geometry of Neuroscience}: neurons and signalling, behavior, the senses, methods and disorders; and Matilde Marcolli's Caltech course on dynamical, topological, geometric and categorical models of the brain, vision and language
        \item[2026]     \tab{}\textit{Genetic Engineering}: notes from MIT 20.205[J]/9.26 (genome editing, CRISPR, and their neuroscience applications)
        \item[ongoing]  \tab{}\textit{Measure Theory}: measure and integration following Folland's \textit{Real Analysis}
        \item[ongoing]  \tab{}\textit{Integrated Mathematics}: an open curriculum from arithmetic through Integrated Maths~I and~II (sets and proof, equations and functions, quadratics, polynomials, sequences and induction, trigonometry, polar coordinates and complex numbers, the real numbers, integral domains)
        \item[ongoing]  \tab{}\textit{Category Theory} and \textit{Biology} (molecular biology, cell biology, genetics, biochemistry)
    \end{tablist}

    \section*{Research Areas and Interests}
    \begin{itemize}
        \item Biostatistical and simulation methods for critical-care medicine: causal inference, pharmacokinetic/pharmacodynamic modeling, and survival optimization
        \item Mathematical modeling of cognition and the mind; mathematical and theoretical neuroscience
        \item Mathematical physics and algebraic foundations: polyadic algebraic structures, polyadic semigroups, and possible extensions to Lie polyadic algebras for differential equations
        \item Philosophy of natural science: Hempel's models of scientific explanation with classical Aristotelian logic replaced by non-classical logics
        \item Education philosophy, politics of education, and political philosophy; post-post modernism
        \item Neurobioengineering and cognitive science
    \end{itemize}

    \section*{Manuscripts, Publications, and Other Writings}
    \subsection*{Books and Textbooks}
    \begin{tablist}
        \item[2025--26] \tab{}Samuel~D.~V.\ Stephens. \textit{Studies in Algebra and Group Theory}. A 519-page textbook written alongside Harvard's Math 55a (fall 2025); revised edition, November 2026. The course instructor agreed to offer it to future Math 55a students as supplementary material. Online at \href{https://www.samuel-stephens.com/notes/math-55a/}{samuel-stephens.com/notes/math-55a}.
    \end{tablist}

    \subsection*{Manuscripts Under Review}
    \begin{tablist}
        \item[2026] \tab{}Samuel~D.~V.\ Stephens\textsuperscript{*}, Tara Westover\textsuperscript{*}, Sahar F.\ Zafar, M.~Brandon Westover. \enquote{CICADAS-II: Simulation-Based Causal Inference and Adaptive Dosing from Sparse Seizure Counts and Bolus Anti-Seizure Medication in the ICU.} Under major revision at the \textit{Journal of Neural Engineering}. \textsuperscript{*}Co-first authors.
    \end{tablist}

    \subsection*{Manuscripts in Preparation}
    \begin{tablist}
        \item[2026--]   \tab{}Samuel~D.~V.\ Stephens. \textit{Studies in Real and Complex Analysis}. Companion textbook for Harvard's Math 55b (spring 2026).
        \item[2026--]   \tab{}Samuel~D.~V.\ Stephens. \enquote{Monotone Arithmetic Progressions in the Gray Code and in Negabinary Enumerations of the Integers.} On Erdős Problems \#195 and \#196.
        \item[2026--]   \tab{}Samuel~D.~V.\ Stephens. \enquote{On Justice: A Manifesto.} On analytic philosophy's failure to provide substance for societal and structural just formation.
        \item[2024--]   \tab{}Samuel~D.~V.\ Stephens, John Cervantes. \enquote{Collaborative Manuscript in Political Philosophy.} A long-term two-person collaboration on the intersections of political theory, educational theory, and ethics, with an extensive chapter outline and drafted sections.
    \end{tablist}

    \subsection*{Musical Arrangements}
    Selected arrangements with publicly available sheet music and accompanying recordings; some works published under the pseudonyms \textit{Czyzew}, \textit{Dr. Grimmmmm}, and \textit{Viktor Alexandre Saint-Yves d'Alveydre}. Catalog of works at \href{https://www.samuel-stephens.com/music}{samuel-stephens.com/music}.
    \begin{itemize}
        \item \textit{Ground} after John Eccles---solo piano arrangement (A-flat minor, 3/4) on a recurring baroque bass figure
        \item \textit{Couplets de folies} after Marin Marais---harp arrangement
        \item \textit{\'Etude d'Ex\'ecution Transcendante No.~1: Grande \'Etude en Fa Mineur} after Henselt---solo piano (under \textit{Czyzew})
        \item \textit{\'Etude d'Ex\'ecution Transcendante No.~2} after Henselt---solo piano (under \textit{Czyzew})
        \item \textit{Armide, LWV 71} after Jean-Baptiste Lully---arrangements of select sections
        \item \textit{Chamber Concerto in D minor, RV 96} after Vivaldi---\href{https://www.youtube.com/watch?v=MfAI0IuN4n8}{recording} and \href{https://musescore.com/user/35512053/scores/6702147}{sheet music} available openly online
    \end{itemize}

    \subsection*{Essays}
    Essays and commentary on philosophy, politics, and mathematics at \href{https://www.samuel-stephens.com/blog}{samuel-stephens.com/blog}.

    \section*{Technical Projects}
    \subsection*{Research and Academic Software}
    \begin{tablist}
        \item[2025--]   \tab{}\textbf{cicadas2} --- Reference implementation supporting the CICADAS-II manuscript and the Stanford research program. Primary developer. \textit{Repository currently private; will be published with the manuscript.}
    \end{tablist}
    \begin{itemize}
        \item Two complete implementations: MATLAB reference (the original) and a verified Python port (NumPy/SciPy), each with a full pytest/MATLAB cross-fixture test suite that locks numerical equivalence
        \item Implements an Extended Kalman Filter on the log-rate state, one-compartment PK with infusion and bolus dosing, Hill/linear PD suppression, five treatment strategies (including adaptive variants), and parametric \textit{g-formula} recovery of causal treatment effects from confounded observational data
        \item Includes the \LaTeX{} manuscript draft, a reproducibility harness that regenerates every reported number, and a figure-generation pipeline in both MATLAB and Python
    \end{itemize}

    \begin{tablist}
        \item[2023--]   \tab{}\textbf{uni} --- Command-line tool for managing coursework, lecture notes, and problem sets in \LaTeX{} and Org-mode. Sole developer; built for personal use and released openly for others. \href{https://github.com/sdvstephens/uni}{github.com/sdvstephens/uni}
    \end{tablist}
    \begin{itemize}
        \item Built an integrated, terminal-based academic workflow as a single Python CLI (\texttt{acad}, exposed through a \texttt{uni} fish-shell wrapper) consolidating course management, lecture and problem-set scaffolding, grade and GPA tracking, and a live status dashboard, in daily personal use across coursework, audits, and tutoring
        \item Generates lecture and problem-set documents in both Org-mode and \LaTeX{} from the command line, with master-file assembly, plain \texttt{pdflatex} via \texttt{latexmk} against a shared preamble, and an Ipe figure-insertion pipeline driven from first-class \texttt{neovim}/vimtex integration
        \item Integrated with Taskwarrior and Timewarrior for deadline-aware, time-tracked course management surfaced through the dashboard
        \item Stores course metadata, grades, and notes as human-readable JSON under Git version control
        \item Published as a FOSS project under MIT---\texttt{pipx install acad-cli}; \href{https://github.com/sdvstephens/acad-cli}{\texttt{sdvstephens/acad-cli}}---inspired by the late Gilles Castel
    \end{itemize}

    \subsection*{Personal Projects}
    \begin{tablist}
        \item[2026--]   \tab{}\textbf{graphite} --- Graphing and 3D calculator for the reMarkable Paper Pro. Sole developer. \textit{Repository currently private; in active development.}
    \end{tablist}
    \begin{itemize}
        \item Renders 2D, parametric, polar, and implicit curves and shaded 3D surfaces, and is designed to insert graphs into the active notebook or document as native, resizable objects
        \item QML (Qt Quick Canvas) frontend for the interface and rendering; Python backend (\texttt{sympy}/\texttt{numpy}/\texttt{scipy}) over a Unix socket for parsing, symbolic computation, sampling, and 3D meshes; delivered to the device through XOVI/AppLoad
    \end{itemize}

    \begin{tablist}
        \item[2024--]   \tab{}\textbf{Better Campus} (formerly BetterCanvas) --- Safari Web Extension for Canvas LMS. Sole developer. \textit{Repository currently private; will be published once stable.}
    \end{tablist}
    \begin{itemize}
        \item Native Safari Web Extension (Xcode/Swift host app plus a Safari extension target) that improves the day-to-day Canvas LMS reading and submission workflow on macOS and iOS/iPadOS
    \end{itemize}

    \begin{tablist}
        \item[2025--26] \tab{}\textbf{rmpp-stabilizer} and \textbf{rmpp-wikipedia} --- Tools for the reMarkable Paper Pro. Sole developer. \textit{Discontinued.}
    \end{tablist}
    \begin{itemize}
        \item \texttt{rmpp-stabilizer}: an \texttt{LD\_PRELOAD} shared library (C++, MIT licensed) that intercepts evdev pen input before \texttt{xochitl} and smooths handwriting in real time, with moving-average, Gaussian-weighted, Krita-style \textit{string pull}, and One Euro (Casiez et al., 2012) filters; Docker-based aarch64 toolchain
        \item \texttt{rmpp-wikipedia}: an offline Wikipedia reader, a Qt/QML interface talking over loopback HTTP to a native \texttt{zimserver} backend built on Kiwix \texttt{kiwix-tools} and \texttt{libzim}
    \end{itemize}

    \section*{Awards and Honors}
    \begin{tablist}
        \item[2025] \tab{}Honors and Dean's Scholar, Wabash College
        \item[2023] \tab{}KEEP-STEM Scholar, Kenyon College (program supported by the National Science Foundation)
        \item[2022] \tab{}AP Capstone Scholar, The College Board
        \item[2021] \tab{}AP Scholar with Distinction and Honors, The College Board
    \end{tablist}

    \section*{Teaching}
    \subsection*{Tutoring}
    \begin{tablist}
        \item[2025--] \tab{}\textit{Private Mathematics Tutor} \\
                            Tutors a Boston exam-school student in precalculus, AP Calculus BC, and physics; writes student guides, worked solutions, and inquiry-based homework.
        \item[2023]   \tab{}\textit{Online Lecturer}, children's mathematics lecture series \\
                            Created a YouTube lecture series (\href{https://www.youtube.com/@Vichaama}{@Vichaama}) with custom curricula and materials.
    \end{tablist}

    \subsection*{Independent Curriculum}
    Open mathematics curriculum for self-paced learners, starting from arithmetic. It takes a structural approach to mathematics, shaped by the shortcomings of the 1960s \textit{New Math}, and a Humboldtian, dialectical view of education. Worksheets, problem sets, and lecture notes are published at \href{https://www.samuel-stephens.com/notes/integrated-maths/}{samuel-stephens.com/notes/integrated-maths}; lectures are scripted in \LaTeX{} and will be recorded for YouTube (\href{https://www.youtube.com/@Vichaama}{@Vichaama}).
    \begin{itemize}
        \item \textit{Basic Mathematics} (in progress)---pre-algebra, arithmetic as structure, algebraic foundations, number theory, geometric reasoning, counting and probability, and functions, laid out lecture by lecture
        \item \textit{Integrated Mathematics I} (notes online)---sets and proof; equations, functions, and graphs; quadratics; polynomials; sequences, series, and induction; with graph theory, combinatorics, and probability planned
        \item \textit{Integrated Mathematics II} (notes online)---trigonometry; parametric equations, polar coordinates, and complex numbers; the real numbers; groups, rings, and fields; with vector spaces and topology planned
        \item \textit{Integrated Mathematics III and IV} (planned)---advanced linear algebra, advanced group theory, real analysis, and topology; then complex analysis, differential equations, Lie groups, geometry, and differential geometry
    \end{itemize}

    %\section*{Service}
   % \subsection*{Board of Directors}
   % \forthcoming{}
    %\section*{Memberships}
    %\forthcoming{}
    % \begin{itemize}
    %     \item North American Computational Linguistics Olympiad (NACLO) --- past competitor
    % \end{itemize}

    \section*{Professional Employment}
    \begin{tablist}
        \item[2025--]   \tab{}Self-employed, Massachusetts \\
                              Private Mathematics Tutor
        \item[2025]     \tab{}1832 Brew at Wabash College, Crawfordsville, Indiana \\
                              Part-Time Barista
        \item[2024]     \tab{}Walgreens, Dallas, Texas \\
                              Sales Representative
        \item[2024]     \tab{}Fossil Group, Inc., Dallas, Texas \\
                              Stylist and Sales Expert
        \item[2023]     \tab{}Nordstrom, E-Bar Coffeehouse, Dallas, Texas \\
                              Summer Barista
    \end{tablist}

    % display today's date as Month Year after a vertical space below the end of the text
    \begin{center}
        \vfill
        Updated \monthyeardate\today
    \end{center}
\end{document}
